\begin{abstract}
\noindent
NEXUS is a drug-discovery pipeline written against NumPy alone: \codeLines{} lines
implementing SMILES parsing and ring perception, molecular fingerprints and
descriptors, reverse-mode automatic differentiation, a directed message-passing
network, Gaussian processes and active learning, split-conformal uncertainty,
persistent homology, discrete Ricci curvature, non-equilibrium free-energy
estimators, a discrete graph diffusion model, and a Gene Ontology layer. It covers
hit identification, lead optimisation, binding-affinity prediction and de novo
generation, and ships as a panel for the Mantis map interface.

The point of the exercise is not that any one component is novel. It is that every
component is checked against something external --- published descriptor tables,
analytically solvable free energies, finite-difference gradients, known topology,
hand-computed hypergeometric tails --- and that the results are reported whichever
way they come out. A D-MPNN reproduces published MoleculeNet numbers, loses to a plain
fingerprint baseline on the smaller datasets, and wins on the largest by beating both
the baselines and the published figure --- the regime crossover made visible rather
than asserted. Conformal intervals attain their
guaranteed coverage, and the adaptive variant that should improve on them helps on
one dataset of three and the textbook construction is worse on all of them. The
generator first scored \emph{below} its own untrained control; that was two bugs, and
fixing them took it to \genModelValidity{} validity against a \genBaselineValidity{}
control, with every remaining failure a single nameable gap. Each of these is reported here as measured,
because a pipeline that only produces good numbers is not a pipeline anyone can use.
\end{abstract}
